\documentclass[11pt,a4paper]{amsart}
\pdfinfoomitdate=1
\pdftrailerid{}
\pdfsuppressptexinfo=15
\usepackage[T1]{fontenc}
\usepackage{lmodern}
\usepackage[margin=27mm]{geometry}
\usepackage{microtype}
\usepackage{amsmath,amssymb,mathtools}
\usepackage{booktabs}
\usepackage{xcolor}
\usepackage[numbers,sort&compress]{natbib}
\usepackage[colorlinks=true,linkcolor=blue!55!black,citecolor=blue!55!black,urlcolor=blue!55!black]{hyperref}
\usepackage[nameinlink,noabbrev]{cleveref}
\raggedbottom

\newtheorem{theorem}{Theorem}[section]
\newtheorem{proposition}[theorem]{Proposition}
\newtheorem{corollary}[theorem]{Corollary}
\newtheorem{lemma}[theorem]{Lemma}
\theoremstyle{definition}
\newtheorem{definition}[theorem]{Definition}
\theoremstyle{remark}
\newtheorem{remark}[theorem]{Remark}

\newcommand{\CP}{\operatorname{CP}}
\newcommand{\Fix}{\operatorname{Fix}}
\newcommand{\Bin}{\operatorname{Bin}}
\newcommand{\Prob}{\mathbb P}

\title[Cocktail-party majority dynamics]{Exact State Partition, Zeta Function, and Critical Window for Majority Dynamics on Cocktail-Party Graphs}
\author{Anonymous}
\date{Internal Stage 2 draft, 27 August 2026}
\subjclass[2020]{37B15, 05C82, 68Q87, 60C05}
\keywords{majority dynamics, cocktail-party graph, threshold network, two-cycle, Artin--Mazur zeta function, critical window}

\begin{document}

\begin{abstract}
Let $\CP_n=K_{2n}\setminus M$ be the cocktail-party graph obtained by
deleting a perfect matching, and apply synchronous strict majority with
inertia at a tie.  We determine the image of every one of its $4^n$ Boolean
states after one step.  The recurrent core consists of $2+2^n$ fixed states
and
\[
 \frac{1}{2}\left\{\binom{2n}{n}-2^n\right\}
\]
genuine two-cycles; every other state reaches consensus in one step.  This
gives exact consensus-basin sizes, every iterate fixed-point count, the
Artin--Mazur zeta function, and the complete symbolic natural extension.  For
i.i.d. Bernoulli initial data we also obtain exact outcome probabilities.  In
the critical window $p=1/2+a/\sqrt n$, the probability of one-consensus tends
to $\Phi(2\sqrt2a)$, while the genuine two-cycle probability is asymptotic to
$e^{-4a^2}/\sqrt{\pi n}$.  An exhaustive implementation of the literal local
rule checks all states through $n=8$ and the first twelve iterate counts.
\end{abstract}

\maketitle

\section{Introduction}

Synchronous Boolean threshold networks with symmetric interaction have no
cycles longer than two, a foundational fact due to Goles and Olivos
\cite{GolesOlivos1980}.  For majority dynamics with the convention that a
vertex retains its state at a tie, Ginosar and Holzman study the same local
rule on infinite graphs \cite{GinosarHolzman2000}.  Quantitative work on
finite majority processes has emphasized convergence time and graph modules
\cite{KaaserEtAl2016}.  Those general results tell us the maximum possible
period, but do not by themselves enumerate the recurrent states or their
basins on a specified graph family.

This note performs that enumeration for the cocktail-party graph.  Its
vertices come in $n$ nonadjacent mate pairs, while every two vertices from
different pairs are adjacent.  A global imbalance forces consensus in one
step.  At exact balance, each mate pair evolves independently: a monochrome
pair flips and a mixed pair freezes.  This elementary dichotomy yields a
complete state partition and, from it, several dynamical quantities that are
usually studied separately.

Our scope is intentionally exact and family-specific.  We use the general
period-two theorem as prior art, and do not treat irreversible conversion
processes, whose multipartite theory has a different update rule
\cite{AdamsEtAl2011}.  The residual contribution is the joint package of
state counts, basin counts, zeta function, natural extension, and Bernoulli
critical-window asymptotics for the symmetric synchronous rule on
$\CP_n$.

\section{The local rule}

Fix $n\geq1$.  Write the vertices of $\CP_n$ as
\[
 V_n=\{1,\bar1,\ldots,n,\bar n\},
\]
where $i$ and $\bar i$ are mates and hence are the only distinct
nonadjacent pair involving either vertex.  A state is
$x=(x_v)_{v\in V_n}\in\{0,1\}^{V_n}$.  Let
\[
 S(x)=\sum_{v\in V_n}x_v.
\]
Every vertex has degree $2n-2$.  The synchronous majority map
$F_n:\{0,1\}^{V_n}\to\{0,1\}^{V_n}$ is
\[
 (F_nx)_v=
 \begin{cases}
 1,&\sum_{w\sim v}x_w>n-1,\\
 0,&\sum_{w\sim v}x_w<n-1,\\
 x_v,&\sum_{w\sim v}x_w=n-1.
 \end{cases}
\]
Thus a tie is inertial.  This convention also makes $F_1$ the identity on
the four states of the edgeless graph $\CP_1$.

\begin{lemma}[one-step trichotomy]\label{lem:trichotomy}
For every $x\in\{0,1\}^{V_n}$:
\begin{enumerate}
\item if $S(x)<n$, then $F_nx=0^{V_n}$;
\item if $S(x)>n$, then $F_nx=1^{V_n}$;
\item if $S(x)=n$, then on each mate pair $\{i,\bar i\}$ the map fixes
the patterns $01$ and $10$ and interchanges $00$ with $11$.
\end{enumerate}
\end{lemma}

\begin{proof}
The number of one-neighbors seen by $v$ is
\[
 N_v(x)=S(x)-x_v-x_{\bar v}.
\]
If $S(x)<n$, then $N_v(x)\leq n-1$.  Equality can occur only when
$S(x)=n-1$ and the mate pair of $v$ is $00$, so inertia still leaves $v$
equal to zero.  Otherwise $N_v(x)<n-1$ and $v$ becomes zero.  This proves
the first assertion; complementation proves the second.

Suppose $S(x)=n$.  A $00$ pair has $N_v(x)=n$ at both vertices and becomes
$11$.  A $11$ pair has $N_v(x)=n-2$ and becomes $00$.  A mixed pair has
$N_v(x)=n-1$ at both vertices, so both coordinates retain their values.
\end{proof}

\section{Complete finite dynamics}

Call a state \emph{pairwise mixed} if each mate pair is either $01$ or
$10$.  There are exactly $2^n$ such states, all of weight $n$.

\begin{theorem}[state partition]\label{thm:partition}
For $n\geq1$, the recurrent core of $F_n$ has:
\begin{align*}
 \#\{\text{fixed states}\}&=2+2^n,\\
 \#\{\text{genuine two-cycles}\}
 &=\frac{1}{2}\left\{\binom{2n}{n}-2^n\right\}.
\end{align*}
The fixed states are the two consensus states and the $2^n$ pairwise-mixed
states.  Each full consensus basin has cardinality
\[
 B_n=\frac{4^n-\binom{2n}{n}}{2}.
\]
Every nonrecurrent state enters the recurrent core after exactly one step,
and $F_n^3=F_n$.
\end{theorem}

\begin{proof}
By \cref{lem:trichotomy}, all states of weight below $n$ enter zero-consensus
and all states of weight above $n$ enter one-consensus.  Complementation is a
bijection between these two sets, while the middle layer has size
$\binom{2n}{n}$.  Hence each consensus basin, including its consensus state,
has size
\[
 \sum_{j=0}^{n-1}\binom{2n}{j}
 =\frac{4^n-\binom{2n}{n}}2.
\]

On the middle layer, pairwise-mixed states are fixed.  Every other state has
at least one monochrome mate pair; all monochrome pairs flip and all mixed
pairs freeze.  A second step therefore returns the original state, but the
first step does not.  Thus the remaining
$\binom{2n}{n}-2^n$ middle-layer states form genuine two-cycles.  The same
description shows that every image is recurrent and proves $F_n^3=F_n$.
\end{proof}

The formula is more informative than the general statement that symmetric
threshold networks have period at most two: it records the complete
functional graph up to the individual leaves attached to the two consensus
states.

\section{Iterate counts, zeta function, and natural extension}

For a map on a finite set, define its Artin--Mazur zeta function by
\[
 \zeta_{F_n}(z)
 =\exp\left(\sum_{k\geq1}\frac{|\Fix(F_n^k)|}{k}z^k\right).
\]

\begin{corollary}[iterate counts and zeta]\label{cor:zeta}
For every $k\geq1$,
\[
 |\Fix(F_n^k)|=
 \begin{cases}
 2+2^n,&k\text{ odd},\\
 2+\binom{2n}{n},&k\text{ even}.
 \end{cases}
\]
Consequently
\[
 \zeta_{F_n}(z)
 =(1-z)^{-(2+2^n)}
  (1-z^2)^{-\frac12\{\binom{2n}{n}-2^n\}}.
\]
\end{corollary}

\begin{proof}
An odd iterate fixes exactly the fixed states.  An even iterate additionally
fixes every point on a two-cycle.  The Euler product of a finite dynamical
system contributes $(1-z^\ell)^{-1}$ for each cycle of length $\ell$, giving
the displayed formula.
\end{proof}

There is also a canonical invertible symbolic system associated with the
noninvertible map.  Its natural extension is
\[
 \widehat X_n
 =\{(x_t)_{t\in\mathbb Z}:F_n(x_t)=x_{t+1}\text{ for every }t\},
 \qquad \widehat\sigma((x_t)_t)=(x_{t+1})_t.
\]

\begin{corollary}[symbolic natural extension]\label{cor:natural}
The projection $(x_t)_t\mapsto x_0$ conjugates
$(\widehat X_n,\widehat\sigma)$ to the restriction of $F_n$ to its recurrent
core.  Hence $\widehat X_n$ is the disjoint union of $2+2^n$ fixed orbits and
$\frac12\{\binom{2n}{n}-2^n\}$ orbits of length two, and its zeta function is
the one in \cref{cor:zeta}.
\end{corollary}

\begin{proof}
In a finite functional graph, a point has prehistories of every length that
fit into a bi-infinite orbit if and only if it lies on a directed cycle.
Theorem~\ref{thm:partition} identifies all such points.  On this recurrent
core $F_n$ is bijective, so the zeroth coordinate uniquely determines the
entire bi-infinite orbit.
\end{proof}

\section{Bernoulli initial data and the critical window}

Let the $2n$ initial coordinates be independent Bernoulli variables of
parameter $p\in[0,1]$.  Write $H_n\sim\Bin(2n,p)$.

\begin{theorem}[exact outcome law]\label{thm:bernoulli}
Under Bernoulli$(p)$ initial data,
\begin{align*}
 \Prob(\text{one-consensus})&=\Prob(H_n>n),\\
 \Prob(\text{zero-consensus})&=\Prob(H_n<n),\\
 \Prob(\text{nonconsensus fixed state})&=[2p(1-p)]^n,\\
 \Prob(\text{genuine two-cycle})
 &=\left\{\binom{2n}{n}-2^n\right\}[p(1-p)]^n.
\end{align*}
These four events partition the probability space.
\end{theorem}

\begin{proof}
The first two identities are \cref{lem:trichotomy}.  Every pairwise-mixed
state contains exactly $n$ ones and has probability $p^n(1-p)^n$; there are
$2^n$ such states.  Every other balanced state lies on a genuine two-cycle,
and there are $\binom{2n}{n}-2^n$ of them.
\end{proof}

At $p=1/2$, the genuine two-cycle probability is
\[
 \frac{\binom{2n}{n}-2^n}{4^n}
 =\frac{1}{\sqrt{\pi n}}(1+O(n^{-1}))-2^{-n}.
\]
The next result resolves the whole $n^{-1/2}$ bias window.

\begin{corollary}[critical window]\label{cor:window}
Fix $a\in\mathbb R$ and put $p_n=1/2+a/\sqrt n$, for all sufficiently large
$n$.  Then
\begin{align*}
 \Prob_{p_n}(\text{one-consensus})&\longrightarrow\Phi(2\sqrt2a),\\
 \Prob_{p_n}(\text{zero-consensus})&\longrightarrow\Phi(-2\sqrt2a),\\
 \Prob_{p_n}(\text{genuine two-cycle})
 &\sim\frac{e^{-4a^2}}{\sqrt{\pi n}}.
\end{align*}
The nonconsensus fixed-state probability is exponentially small.
\end{corollary}

\begin{proof}
For $H_n\sim\Bin(2n,p_n)$,
\[
 \frac{n-2np_n}{\sqrt{2np_n(1-p_n)}}\longrightarrow-2\sqrt2a.
\]
The central limit theorem gives the first two limits; the atom at $n$ is
$O(n^{-1/2})$ and does not affect them.  Stirling's formula and
\[
 [p_n(1-p_n)]^n
 =4^{-n}(1-4a^2/n)^n
 \sim4^{-n}e^{-4a^2}
\]
give the two-cycle asymptotic.  Finally,
$[2p_n(1-p_n)]^n=2^{-n}(1-4a^2/n)^n$.
\end{proof}

\section{Deterministic control and claim boundary}

The companion script constructs the deleted-matching neighbor list and
evaluates the literal local majority rule; it does not use the trichotomy as
an update shortcut.  It exhausts all $4^n$ states for $1\leq n\leq8$, checks
$F_n^3=F_n$, the state and basin formulas, and $|\Fix(F_n^k)|$ for
$1\leq k\leq12$.  This is a finite regression control, not a substitute for
the proofs.

General period-two behavior is explicitly attributed to
\citet{GolesOlivos1980}, and the tie convention to the established majority
literature \cite{GinosarHolzman2000}.  A bounded search through 27 August
2026 found no exact collision with the combined cocktail-party formulas in
\cref{thm:partition,cor:zeta,cor:natural,thm:bernoulli,cor:window}; this is a
search record, not a worldwide priority claim.  Public release and venue
selection remain outside the scope of this internal draft.

\bibliographystyle{plainnat}
\bibliography{references}

\end{document}
